\documentclass{article}
\usepackage{amsmath,amssymb,amsthm}

\theoremstyle{plain}
\newtheorem{theorem}{Theorem}[section]
\newtheorem{lemma}[theorem]{Lemma}
\newtheorem{proposition}[theorem]{Proposition}
\newtheorem{corollary}[theorem]{Corollary}
\newtheorem*{conjecture}{Conjecture}

\theoremstyle{definition}
\newtheorem{definition}[theorem]{Definition}
\newtheorem{example}{Example}
\newtheorem*{exercise}{Exercise}

\theoremstyle{remark}
\newtheorem{remark}[theorem]{Remark}
\newtheorem*{note}{Note}

\begin{document}

\section{Orders and lattices}

A partially ordered set is the basic object of this section. We begin with
the definitions and then prove a few elementary facts.

\begin{definition}
A \emph{partial order} on a set $P$ is a relation $\le$ that is reflexive,
antisymmetric and transitive. The pair $(P,\le)$ is a \emph{poset}.
\end{definition}

\begin{definition}[Upper bound]
Let $A \subseteq P$. An element $u \in P$ is an \emph{upper bound} of $A$ if
$a \le u$ for every $a \in A$. The least upper bound, when it exists, is
written $\sup A$.
\end{definition}

\begin{example}
The power set $\mathcal{P}(X)$ ordered by inclusion is a poset in which every
family $\mathcal{F}$ has the supremum $\bigcup \mathcal{F}$.
\end{example}

\begin{lemma}
A subset of a poset has at most one least upper bound.
\end{lemma}

\begin{proof}
If $u$ and $v$ are both least upper bounds of $A$, then $u \le v$ and
$v \le u$, so $u = v$ by antisymmetry.
\end{proof}

\begin{theorem}[Knaster--Tarski]
Let $L$ be a complete lattice and let $f \colon L \to L$ be order preserving.
Then the set of fixed points of $f$ is itself a complete lattice; in
particular $f$ has a least fixed point.
\end{theorem}

\begin{remark}
Completeness cannot be dropped: the map $n \mapsto n+1$ on $\mathbb{N}$
preserves order and has no fixed point.
\end{remark}

\begin{corollary}
Every order preserving map on a finite lattice has a fixed point.
\end{corollary}

\section{Chains}

\begin{proposition}\label{prop:chain}
Every finite nonempty chain has a greatest element.
\end{proposition}

\begin{proof}
Induct on the size $n$ of the chain. For $n = 1$ there is nothing to show.
Otherwise remove any element $x$, take the greatest element $m$ of the rest,
and compare $x$ with $m$.
\end{proof}

\begin{example}[Integers]
The chain $(\mathbb{Z},\le)$ has no greatest element, so finiteness in
Proposition~\ref{prop:chain} is needed.
\end{example}

\begin{conjecture}
Every sufficiently regular poset of width $w$ can be covered by $w$ chains
that are computed greedily.
\end{conjecture}

\begin{exercise}
Show that the divisibility order on $\{1,2,\dots,12\}$ is not a chain.
\end{exercise}

\begin{note}
Starred environments are unnumbered and do not advance any counter.
\end{note}

\begin{theorem}
Every poset can be extended to a total order on the same set.
\end{theorem}

\end{document}
